% Hodina 63: rozpis 1. ročníku, řádek 63.
% Zdroje: ucebnice/02_Funkce_a_rovnice/Linearni_funkce_a_rovnice/Lineární funkce.tex,
% oddíl Lineární funkce; předchozí Funkce.tex; lokální Realisticky.cz
% 020109_Linearni_funkce_II/020109_Lekce.pdf (zavedení tvaru ax+b).
% Úlohy a graf vlastní tvorba.
\documentclass[11pt,a4paper]{article}
% Spolecny styl pro ucitelske pripravy MAT (Work/MAT)
% Verze 1.1, 2026-10-23
\NeedsTeXFormat{LaTeX2e}
\ProvidesPackage{priprava}[2026/10/23 v1.1 Sablona ucitelskych priprav MAT]

% Zakladni typografie a cestina
\RequirePackage[T1]{fontenc}
\RequirePackage[utf8]{inputenc}
\RequirePackage[czech]{babel}
\RequirePackage{lmodern}
\RequirePackage{microtype}
\RequirePackage{amsmath,amssymb}
\RequirePackage{xcolor}
\RequirePackage{enumitem}
\RequirePackage[most]{tcolorbox}
\RequirePackage{tabularx,booktabs,array}
\RequirePackage{fancyhdr}
\RequirePackage{needspace}
\RequirePackage{geometry}
\RequirePackage[hidelinks]{hyperref}

\geometry{a4paper,top=17mm,bottom=18mm,left=23mm,right=23mm}
\definecolor{notegray}{gray}{0.39}
\setlength{\parindent}{0pt}
\setlength{\parskip}{2.5pt plus 1pt minus 0.5pt}
\setlist[itemize]{topsep=2pt,itemsep=1pt,parsep=0pt,leftmargin=1.45em}
\setlist[enumerate]{topsep=2pt,itemsep=2pt,parsep=0pt,leftmargin=1.65em}
\setcounter{secnumdepth}{0}
\renewcommand{\arraystretch}{1.12}

% Zahlavi: vlevo rocnik | globalni tema, vpravo cislo: nazev hodiny
% Zapati zustava zcela prazdne.
\pagestyle{fancy}
\fancyhf{}
\renewcommand{\headrulewidth}{0.35pt}
\renewcommand{\footrulewidth}{0pt}
\setlength{\headheight}{14pt}
\newcommand{\nastavPripravu}[4]{%
  \fancyhead[L]{\small #1\enspace\textbar\enspace #2}%
  \fancyhead[R]{\small #3: #4}%
}

% Minimalisticke cernobile boxy. Kazda definice zvlast.
\tcbset{commonlesson/.style={enhanced,breakable=false,colback=white,colframe=black,
  colbacktitle=white,coltitle=black,boxrule=0.45pt,
  arc=0.8mm,outer arc=0.8mm,left=3mm,right=3mm,top=1.6mm,bottom=1.6mm,
  toptitle=1mm,bottomtitle=0.7mm,fonttitle=\bfseries,
  attach boxed title to top left={xshift=2.5mm,yshift=-1.15mm},
  boxed title style={colback=white,colframe=white,boxrule=0pt,left=1mm,right=1mm}}}
\newtcolorbox{zadaniBox}[1][]{commonlesson,title={Zadání},#1}
\newtcolorbox{prikladBox}[1][]{commonlesson,title={Příklad},#1}
\newtcolorbox{definiceBox}[1][]{commonlesson,title={Definice},#1}

% Spolecny box pro zadani a jeho reseni, bez mezery mezi castmi.
% Pouziti:
%   \begin{prikladReseniBox}
%     Zadani ulohy.
%     \tcblower
%     \begin{align*} ... \end{align*}
%   \end{prikladReseniBox}
% Spodni cast ma jemny sedy podklad a delici linku; lze v ni sazet
% vice radku matematiky. Obe casti maji jen jeden spolecny ramecek.
\tcbset{solutionlesson/.style={
  commonlesson,
  enhanced jigsaw,
  breakable,
  colbacklower=black!3,
  segmentation style={solid,draw=black!25,line width=0.35pt},
  fontlower=\small\color{notegray},
  before lower={\noindent\textbf{Řešení:}\enspace},
  leftlower=3mm,rightlower=3mm
}}
\newtcolorbox{prikladReseniBox}[1][]{solutionlesson,title={Příklad},#1}
\newtcolorbox{zadaniReseniBox}[1][]{solutionlesson,title={Zadání},#1}

% Pomucky pro autora pripravy: sedive poznamky; \reseni pro starsi pripravy
\newcommand{\poznamka}[1]{\par\begingroup\small\color{notegray}\itshape
  \textnormal{\textbf{Poznámka pro učitele.}}\enspace #1\par\endgroup}
\newcommand{\reseni}[1]{\par\begingroup\small\color{notegray}
  \textbf{Očekávané řešení:}\enspace #1\par\endgroup}
\newcommand{\blok}[1]{\needspace{5\baselineskip}\subsection*{#1}}
\newcommand{\vyrok}[1]{\uv{#1}}

\usepackage{tikz}
\nastavPripravu{1. ročník}{Funkce}{63}{Lineární funkce a její graf}
\begin{document}
\blok{Motivace a návaznost}
Po obecných vlastnostech funkcí se zaměříme na funkce, jejichž hodnota se při stejném přírůstku argumentu mění vždy o stejnou částku. Příklad $f(x)=2x+1$ dává pro $x=0,1,2$ hodnoty $1,3,5$: pokaždé přibyly $2$. Grafem je přímka. Podrobný význam obou parametrů vyložíme v hodině 64.
\begin{definiceBox}
\textbf{Lineární funkce} na $\mathbb R$ je funkce $f(x)=ax+b$, kde $a,b\in\mathbb R$ jsou pevné koeficienty. Její graf je přímka; při $a=0$ je vodorovná. Je-li stejný předpis omezen na podmnožinu $\mathbb R$, grafem je jen příslušná část přímky.
\end{definiceBox}
\begin{prikladReseniBox}
Pro $f(x)=2x+1$ sestavte tabulku hodnot při $x=-1,0,1,2$. Které body patří grafu?
\tcblower
\[
\begin{array}{c|rrrr}x&-1&0&1&2\\\hline f(x)&-1&1&3&5\end{array}
\]
Dosazení dává $f(-1)=-1$, $f(0)=1$, $f(1)=3$, $f(2)=5$. Graf obsahuje body $(-1;-1),(0;1),(1;3),(2;5)$, ležící na jedné přímce.
\end{prikladReseniBox}
\begin{center}
\begin{tikzpicture}[x=.75cm,y=.52cm,>=stealth]
\draw[step=1,gray!25,very thin] (-3,-3) grid (3,5);
\draw[->] (-3.4,0)--(3.5,0) node[right] {$x$};\draw[->] (0,-3.4)--(0,5.4) node[above] {$y$};
\foreach \x in {-2,-1,1,2}{\draw (\x,.08)--(\x,-.08) node[below] {\small $\x$};}
\foreach \y in {-1,1,3,5}{\draw (.08,\y)--(-.08,\y) node[left] {\small $\y$};}
\draw[very thick] (-2,-3)--(2,5);\foreach \x/\y in {-1/-1,0/1,1/3,2/5}{\fill (\x,\y) circle (2.3pt);}
\node[below right=3pt] at (-1,-1) {\small $(-1;-1)$};\node[above right] at (1,3) {\small $(1;3)$};
\end{tikzpicture}
\end{center}
\blok{Předpis, argument a graf}
V zápisu $f(x)=ax+b$ má $x$ roli argumentu, $f(x)$ funkční hodnoty. Pro každé reálné $x$ umíme hodnotu vypočítat, proto přirozený definiční obor tohoto předpisu je $\mathbb R$. Přímku pro určení tvaru stačí načrtnout pomocí dvou různých bodů; při ruční konstrukci můžeme třetím bodem ověřit výpočet či měřítko.
\begin{prikladReseniBox}
Jsou $p(x)=-3x+2$, $q(x)=x^2+1$ a $r(x)=\frac{x+1}{2}$ lineární funkce na $\mathbb R$? U každé lineární určete hodnoty v $0$ a $2$.
\tcblower
$p$ je tvaru $ax+b$ s $a=-3,b=2$; $p(0)=2$, $p(2)=-4$. $q$ lineární není: obsahuje $x^2$ a její graf není přímka. $r(x)=\frac12x+\frac12$ lineární je; $r(0)=\frac12$, $r(2)=\frac32$. Zlomek předpisu nevadí, pokud je jmenovatel konstantní nenulový.
\end{prikladReseniBox}
\begin{prikladReseniBox}
Najděte předpis lineární funkce, o níž víte, že $f(0)=3$ a $f(1)=5$.
\tcblower
Píšme $f(x)=ax+b$. Z $f(0)=b=3$ dostaneme $b=3$. Dále $f(1)=a+3=5$, tedy $a=2$. Hledaná funkce je $f(x)=2x+3$. Kontrola: $f(0)=3$, $f(1)=5$.
\end{prikladReseniBox}
\poznamka{Nestačí říci, že graf „vypadá rovně“. Přímkovost plyne z tvaru $ax+b$; třetí tabulkový bod kontroluje kresbu. Podrobná konstrukce z tabulky přijde v hodině 65.}
\blok{Samostatné procvičení}
\begin{zadaniReseniBox}
\begin{enumerate}[label=\alph*)]
\item Pro $g(x)=-x+4$ dopočtěte $g(-2),g(0),g(3)$ a napište odpovídající body.
\item Zjistěte, zda body $A(1;4)$ a $B(2;6)$ leží na grafu $h(x)=2x+2$.
\item Určete předpis lineární funkce s $u(0)=-1$ a $u(1)=2$.
\end{enumerate}
\tcblower
\begin{enumerate}[label=\alph*)]
\item $g(-2)=6$, $g(0)=4$, $g(3)=1$; body $(-2;6),(0;4),(3;1)$.
\item $h(1)=4$ a $h(2)=6$, takže oba body leží na grafu.
\item $u(x)=ax+b$, $b=u(0)=-1$. Rovnice $a-1=2$ dává $a=3$, tedy $u(x)=3x-1$. Kontrola $u(1)=2$.
\end{enumerate}
\end{zadaniReseniBox}
\begin{zadaniReseniBox}[title={Diagnostika: přímá úměrnost a rovnoměrná změna}]
\begin{enumerate}[label=\alph*)]
\item Rozhodněte, zda $y=3x$ a $y=3x+2$ procházejí počátkem. Jak se v obou případech změní hodnota při zvýšení $x$ o $1$?
\item Tabulka obsahuje hodnoty $q(0)=1$, $q(1)=2$, $q(2)=5$. Může je všechny popsat jediná lineární funkce na $\mathbb R$?
\end{enumerate}
\tcblower
\begin{enumerate}[label=\alph*)]
\item $y=3x$ prochází $(0;0)$ a jde o přímou úměrnost. Funkce $y=3x+2$ má při $x=0$ hodnotu $2$, proto počátkem neprochází. V obou případech je změna při kroku $1$ rovna $3$.
\item Nemůže. Změny hodnot při stejných krocích argumentu jsou $2-1=1$ a $5-2=3$, ale u lineární funkce musejí být stejné. Tím také konkrétně zdůvodníme, proč předpis $x^2+1$ není lineární na $\mathbb R$.
\end{enumerate}
\end{zadaniReseniBox}
\blok{Volitelné rozšíření a rezerva}
\begin{zadaniReseniBox}
Je funkce $v(x)=\frac{2x^2-2}{x-1}$ lineární funkcí na celé $\mathbb R$? Porovnejte ji s $w(x)=2x+2$ na $\mathbb R$.
\tcblower
Čitatel je $2(x-1)(x+1)$, takže pro $x\ne1$ je $v(x)=2x+2$. Původní funkce v $x=1$ není definována; její graf je přímka s vynechaným bodem $(1;4)$. Proto $v$ není lineární funkcí na celé $\mathbb R$ podle dnešní definice. Funkce $w$ je zadána i v $x=1$ a její graf obsahuje bod $(1;4)$.
\end{zadaniReseniBox}
\blok{Domácí práce}
\begin{zadaniReseniBox}
\begin{enumerate}[label=\alph*)]
\item Pro $f(x)=4-2x$ sestavte tabulku při $x=-1,0,1,2$.
\item Najděte předpis funkce $g(x)=ax+b$ s $g(0)=2$ a $g(2)=8$.
\item Rozhodněte, zda $A(-1;5)$ leží na grafu $h(x)=-2x+3$.
\end{enumerate}
\tcblower
\begin{enumerate}[label=\alph*)]
\item Hodnoty jsou $6,4,2,0$; kontrola dosazením: například $f(2)=4-4=0$.
\item $b=2$, $2a+2=8$, tedy $a=3$ a $g(x)=3x+2$. Kontrola $g(2)=8$.
\item $h(-1)=2+3=5$, proto bod do grafu patří.
\end{enumerate}
\end{zadaniReseniBox}
\poznamka{Jádro 45 min: motivace 4, definice 7, tabulka a graf 12, rozpoznání předpisů 7, samostatné procvičení 8, kontrola 5, organizace 2 min. Rezerva 5–8 min, domácí práce 10–15 min. V další hodině vysvětlíme, jak koeficienty mění sklon a polohu přímky.}
\end{document}
